\section{Discussion}
\label{sec:discussion}

\subsection{What the Acquisition Result Means}

The clearest result is not that explanation guidance replaces uncertainty. It is
that directional attribution balance adds a small, repeatable improvement to the
learning trajectory over random acquisition and a smaller improvement over
uncertainty in four of five outer folds. $H$ distinguishes cases in which the
same output margin arises from opposing signed masses. Such cases can expose a
model boundary supported by competing PE attributes, making them plausible
queries for label acquisition.

The algebraic identity in (\ref{eq:h-identity}) also limits this interpretation.
$H$ is determined by the raw-score contribution and the $L_1$ attribution norm.
Likewise, $V$ is determined by $L_2$ magnitude and the mean signed contribution.
A direct margin-plus-norm score may recover part of the measured gain. Until
those controls are run, the result supports the tested summaries but does not
show that entropy and variance are uniquely necessary.

No policy dominates every objective. $U+H$ leads AULC, $U+V$ reaches the target
with the fewest labels, $\xfull$ performs between them, and Uncertainty remains
strong in restricted-FPR metrics. This is scientifically more informative than a
single winner: acquisition policy should be selected against an explicit
objective and operating region. A system that values early target attainment may
prefer $U+V$, whereas one that values average ranking quality throughout model
development may prefer $U+H$.

The target-based label result must also include its accounting boundary. The
numbers count labels admitted to model training. Validation labels, the full-pool
reference, dataset labeling, and analyst adjudication are additional resources.
A 92-label reduction does not imply a measured reduction in analyst hours.

\subsection{Batch Diversity and Robust Scaling}

The top-100 selection rule can query redundant samples. This risk is material
because the batch equals half of the 200-sample initial set. Uncertainty plus
k-center selection would test whether the observed gain comes from SHAP signals
or whether coverage of the same uncertain region is sufficient. A calibrated
committee would test whether ensemble disagreement explains the benefit.

Min--max normalization is also pool dependent. One extreme $H$ or $V$ value can
compress the remaining component scores and change the effective mixture.
Weight sensitivity, percentile/rank normalization, and robust quantile scaling
should be frozen and evaluated before final submission. These studies must use
validation or a nested design; selecting the best transformation from test AULC
would invalidate the outer comparison.

\subsection{What Controlled Deletion Establishes}

Top-SHAP replacement consistently damages fitted-model ranking more than
uniform random replacement at the same coordinate count, and the result holds
for mean and median baselines. This supports a relative model-reliance claim:
TreeSHAP ranks coordinates that influence these fitted LightGBM models under the
defined intervention.

Several alternatives remain. Tree ensembles use only a subset of the 2,381
coordinates; uniform random masks will often select inactive features.
Top-ranked coordinates can also lie farther from their replacement baselines,
creating a larger perturbation magnitude and greater distribution shift. A
stronger suite would select random controls from model-used features, match
absolute standardized displacement, replace semantically coupled feature
groups, and compare TreeSHAP with gain/permutation rankings.

An attribution-column permutation is a particularly useful negative control. A
common permutation preserves attribution sums, cosine similarities, centroid
distances, and norms, so additivity, EC, and CF can all remain unchanged while
feature identities are wrong. A valid identity-sensitive pipeline should fail
that control. Until it is included, EC and CF are geometry descriptions rather
than explanation certificates.

\subsection{Stability and Model Maintenance}

Moderate within-fold agreement suggests that global attribution summaries are
not dominated by model randomness. The larger reduction across folds shows that
the learned explanation depends materially on the sampled data. This matters for
model maintenance: a global top-feature list should be versioned with the model
and training cohort instead of presented as a timeless malware signature.

The active-learning use of SHAP is local and round specific, so some adaptation
of the acquisition ranking is expected. Nevertheless, a deployment study should
track rank drift, query composition, and class yield across releases. Sharp
changes can signal data drift, feature-extractor changes, or instability that
requires investigation before analyst queries are trusted.

\subsection{Relationship to Multi-Stage Fusion}

The broader thesis studies conditional static/dynamic analysis. A separate
199-sample pilot reported complementary static and API-sequence scores and found
that proportional fusion had the highest mean AUC/AP among seven configurations.
The PI and PID variants produced identical final decisions; the derivative term
added no measurable benefit at 5--20 event prefixes, and the PID-versus-fixed AUC
interval crossed zero. Those results motivate evidence-driven simplification,
not a full PID superiority claim.

That pilot uses MalBehavD-V1, logistic scorers, and event-indexed API prefixes;
the present AL/Fidelity study uses static EMBER vectors and LightGBM. Their
metrics cannot be combined into an end-to-end system result. Multi-stage
integration requires one paired cohort, a gate with final decisions for all
traffic, fixed/learned fusion baselines, and a factorial acquisition--fusion
study. We therefore keep fusion outside the contributions and abstract of this
paper.

\subsection{Practical Implications}

The most immediate use is a model-development workflow. Train a static PE model,
score a candidate pool, apply a policy chosen for the development objective,
present the selected rows to an authorized analyst, and retain query/model
lineage. Controlled deletion and stability checks can run as release diagnostics
for candidate models. They do not authorize automated remediation or substitute
for malware-analysis containment and labeling governance.

The query-computation overhead is modest in the measured small pool, but scaling
TreeSHAP to large corpora can become material. A production design may prefilter
by uncertainty and calculate explanations only for a candidate band, or cache
tree contributions for repeated policy comparisons. Such engineering changes
must be re-evaluated because they change the acquisition population.

